\clearpage
\VerdierPhysicalPageAnchor{0195}%
\begingroup\small\setlength{\parskip}{0pt}\setlength{\emergencystretch}{1.5em}
\setlength{\abovedisplayskip}{4pt plus 2pt minus 2pt}%
\setlength{\belowdisplayskip}{4pt plus 2pt minus 2pt}%

\VerdierSemanticAnchor{ch3d-2-3-1-proof-cont}
où les $X_n$ sont des objets de $\operatorname{comp}^{-}(\mathcal A)$,
les $u_n$ des morphismes élémentaires de hauteur $-n$ et un diagramme
commutatif à homotopie près :

\VerdierSemanticAnchor{ch3d-2-3-1-limits-square}
\begin{center}
\begin{tikzcd}[column sep=5.5em,row sep=3.0em]
\displaystyle\varprojlim_n X_n \arrow[r] \arrow[d]
  & X \arrow[d,"m"] \\
\displaystyle\varinjlim_n X_n \arrow[r,"\sim"'] & Y
\end{tikzcd}
\end{center}

\VerdierSemanticAnchor{ch3d-2-3-1-normalization}
Pour démontrer la proposition, on peut donc supposer que
$X=\varprojlim_n X_n$ et que $X_n=Y$, $u_n=\operatorname{id}$, pour $n$
assez petit. Nous allons construire par récurrence sur $n$ :

\VerdierSemanticAnchor{ch3d-2-3-1-induction-data}
\noindent a) Une suite de morphismes élémentaires :
\[
v_n\colon L_n^\bullet\longrightarrow L_{n-1}^\bullet,
\qquad n\in\mathbf Z,\qquad
v_n\text{ élémentaire de hauteur }-n.
\]

\noindent b) Une suite de morphismes :
\[
w_n\colon L_n^\bullet\longrightarrow X_n
\]
tels que :

\VerdierSemanticAnchor{ch3d-2-3-1-induction-condition-1}
\noindent 1) Pour tout $n\in\mathbf Z$ les composants de degré $i$,
$i\geq -n$, de $L_n^\bullet$ soient des éléments de $M$ et les composants de
degré $i$, $i<-n$, soient des éléments de $M^\wedge$.

\VerdierSemanticAnchor{ch3d-2-3-1-induction-square}
\noindent 2) Les diagrammes ci-après soient commutatifs :
\begin{center}
\begin{tikzcd}[column sep=5.0em,row sep=2.8em]
L_n^\bullet \arrow[r,"w_n"] \arrow[d,"v_n"']
  & X_n \arrow[d,"u_n"] \\
L_{n-1}^\bullet \arrow[r,"w_{n-1}"'] & X_{n-1}
\end{tikzcd}
\end{center}
pour tout $n\in\mathbf Z$.

\VerdierSemanticAnchor{ch3d-2-3-1-induction-condition-3}
\noindent 3) L'objet $L_n^\bullet$ soit égal à $Y$ et les morphismes $v_n$
et $w_n$ soient l'identité pour $n$ assez petit.
\endgroup
